%% =====================================================================
%%  Chapter 3: Related Work
%%  Organize the related studies and explain how this work differs from them.
%%  End the chapter with the gap in existing work that motivates your method.
%% =====================================================================
\chapter{Related Work}
\label{chap:related-work}

This chapter reviews studies related to this thesis and clarifies its position.
\begin{writingnote}
Rather than summarizing papers one by one, group them by viewpoint (e.g., the problem setting or the idea of the method),
give each group its own section, and discuss the ideas, strengths and limitations of each group.
For each work, state what it studied, how and what it found, and what it shares with this study and how it differs.
Point out the limitations of existing work fairly and with evidence, taking into account the aims and assumptions of each study.
Read every work you cite yourself; avoid citing a work only on the basis of how another paper describes it.
Use the present perfect for general statements about the literature (``Many methods have been proposed for \ldots''),
the past tense for what a specific study did (``Dantzig proposed \ldots''),
and the present tense for facts that are still accepted~\cite{glasmandeal2010science,morleyphrasebank}.
\end{writingnote}

\section{Studies on XX}
\label{sec:related-a}

%% Cite with \cite{key}. The key is the name in the first line of each entry in references.bib.
%% Terms and symbols defined in the previous chapter can be used here (e.g., \cref{sec:problem}).
Many methods have been studied for the problem described in \cref{sec:problem}.
For example, Dantzig~\cite{dantzig1957discrete} \ldots.
Kellerer et al.~\cite{kellerer2004knapsack} \ldots.

\section{Studies on YY}
\label{sec:related-b}

%% Japanese works are cited in the same way.
Nishi~\cite{nishi2020} \ldots. Nishi et al.~\cite{nishi2011} \ldots.
%% When you cite several works at once, the numbers are sorted and compressed, e.g., [1-3].
Research on \ldots{} is growing~\cite{nishi2011,nishi2020,dantzig1957discrete}.

\section{Position of This Study}
\label{sec:position}

\begin{writingnote}
Based on the review in the previous sections, state clearly the problem that existing work has not solved,
and explain how this study addresses it, which leads to the Proposed Method chapter.
The problem stated here must match the one stated in the Introduction.
Choose the viewpoints of the comparison because they matter for the problem, not because they favor this study.
Before writing that no study has addressed something, search the literature thoroughly,
and limit the claim, for example, with ``to the best of our knowledge''.
\end{writingnote}

%% A table that compares existing work with yours from several viewpoints is easy to read.
\Cref{tab:comparison} compares related work with this study.
Unlike existing work, which does not consider \ldots, this study considers \ldots.

\begin{table}[tbp]
  \centering
  \caption[Comparison of related work and this study (illustrative)]{Comparison of related work and this study (formatting example; the marks are illustrative and do not assess the cited works)}
  \label{tab:comparison}
  \begin{tabular}{lccc}
    \toprule
    & Viewpoint A & Viewpoint B & Viewpoint C \\
    \midrule
    \cite{dantzig1957discrete} & \checkmark & & \\
    \cite{nishi2011}           & \checkmark & \checkmark & \\
    This study                 & \checkmark & \checkmark & \checkmark \\
    \bottomrule
  \end{tabular}
\end{table}
